\section{Discussion}

\subsection{Weak Relevance, Identification, and Inferential Informativeness}

The common deterioration across Oracle-GMM, Full-GMM, and DML-IVQR-BC is consistent with reduced identifying information. In IVQR, identification depends on whether instrument-induced treatment variation changes the structural quantile moments enough to separate $\alpha_0(\tau)$ from candidate values. When relevance is high, moving $a$ away from the truth tends to move the population moment away from zero, and $W_N(a)$ can reveal that disagreement. As relevance declines, true and false candidates generate more similar moment behavior, though an individual sample profile need not be literally flat. The same mechanism explains three findings: point-estimation error rises, more candidates remain in the inverted accepted set, and false candidates are rejected less often. The design parameter $\kappa$ creates this controlled change in relevance; it is not itself a formal measure of IVQR identification.

Oracle-GMM shows that high-dimensional nuisance estimation is not sufficient to explain the pattern. It knows the designated oracle control set but still experiences increasing error, expanding accepted sets, and declining power. Estimator-specific nuisance and residualization details prevent a perfect decomposition, yet the component shared by all three procedures is weaker information from the excluded instruments. Increasing the sample from $500$ to $1000$ generally improves accuracy and discrimination, but does not eliminate the differences across $\kappa$. Quantile-specific rates and asymmetries also remain, consistent with the density-weighted nature of IVQR information, although the simulation does not construct a formal quantile-specific identification statistic.

The domain diagnostic reinforces this interpretation while limiting the claim appropriately. Expansion adds little accepted measure at $\kappa=1$, substantial measure in some profiles at $\kappa=0.25$, and continuing measure through the largest investigated domain at $\kappa=0.10$. The weakest designs therefore provide little numerical localization within the investigated ranges. They are not thereby proven globally unbounded or formally unidentified.

Robustness and informativeness answer different questions. The candidate-specific procedure of \textcite{chernozhukov_hansen_2008} avoids relying on an approximately normal, strongly identified point estimator. Its limiting null result can remain valid under weak, partial, or failed identification, but it does not promise narrow confidence regions or high power. A wide or disconnected region can correctly report that the maintained model and data exclude few values. Finite-sample calibration is separate again. The preserved procedure compares $W_N(a)$ with an asymptotic $\chi^2_2$ cutoff, and Chapter~5 shows that its finite-sample null distribution can differ substantially from this reference. Direct coverage and the audit rule out grid artifacts, while Oracle undercoverage shows that the issue need not originate in high-dimensional selection. The evidence concerns approximation quality at the simulated sample sizes, not the validity of the limiting theorem. A procedure may consequently be asymptotically robust but uninformative, poorly calibrated in finite samples, or both.

\subsection{What DML Can and Cannot Address}

DML-IVQR targets estimation of high-dimensional nuisance parameters surrounding the low-dimensional structural coefficient. Quantile Lasso regularizes the candidate-specific outcome nuisance, and regularized weighted regressions estimate instrument residualization. Neyman orthogonality reduces the first-order effect of sufficiently small nuisance-estimation errors on the population moment \parencite{chernozhukov_et_al_2018,chen_huang_tien_2021}. This local property neither eliminates nuisance error nor guarantees accurate finite-sample chi-square calibration. More fundamentally, it does not change the variation supplied by the instruments.

The comparison with Full-GMM indicates why nuisance treatment can matter. Full uses all controls without $\ell_1$ regularization, while DML often has lower RMSE, smaller accepted sets, and higher power. The ordering varies across quantiles, relevance settings, and metrics, so DML is not uniformly preferable. Nor can the differences be attributed entirely to regularization: DML and Full also differ in their nuisance and residualization procedures, including intercept treatment. The comparison supports conclusions about the implemented estimators, not a causal decomposition of their components.

Oracle-GMM helps delimit this explanation. By avoiding estimation of the full high-dimensional nuisance model, it removes the central problem that DML is designed to manage, yet it deteriorates with $\kappa$ alongside the other procedures. Orthogonality changes the influence of nuisance-estimation error; it cannot turn weakly informative structural moments into strongly informative ones. The common low-relevance deterioration is therefore not evidence against regularization. It demonstrates that nuisance robustness and instrument relevance are separate requirements.

The DML label here refers to the preserved DML-IVQR-BC implementation. It uses the pivotal penalty, fitted-normal density proxy, and same-sample score evaluation described in Section~3.3; it does not use cross-fitting. These choices define the comparison without implying that their alternatives would improve performance. The results consequently support a bounded conclusion: handling a large nuisance model can improve finite-sample behavior relative to the unregularized full-control procedure, but no nuisance method can restore structural information removed from the instruments.

\subsection{Limitations and Implications}

The experiment is based on one high-dimensional IVQR DGP derived from \textcite{chen_huang_tien_2021}. Holding its treatment distribution, endogeneity channel, outcome equation, structural quantile effect, controls, and instruments fixed across $\kappa$ strengthens the internal interpretation, but limits generalization to other disturbance distributions, degrees of endogeneity, instrument dimensions, sparsity patterns, or structural functions. The two sample sizes do not describe very small samples, larger-sample convergence, or rates. Five quantiles and six power offsets reveal heterogeneity but do not map the continuous quantile or power functions. Monte Carlo summaries based on 500 replications retain simulation uncertainty, which the coverage audit treats explicitly.

The evidence also concerns one same-sample DML implementation rather than all possible nuisance learners or cross-fitting schemes. The three estimators differ in more than regularization, so their comparison does not isolate any one component. Inference uses the preserved asymptotic chi-square cutoff; bootstrap corrections and the exact conditional procedure of \textcite{chernozhukov_hansen_jansson_2009} are not evaluated. The accepted-set measure is computed on a finite grid, and even the expanded domains cannot establish global boundedness or unboundedness. Finally, the thesis contains no empirical application, so it does not measure the frequency or importance of these patterns in a particular economic setting.

These limitations also shape empirical practice. A conventional first-stage diagnostic cannot by itself establish quantile-specific identification, and nominal coverage says little about how many economically distinct effects remain plausible. Reporting the criterion or accepted region over a transparent candidate domain, together with power or sensitivity calculations where feasible, makes the information content visible. When calibration is doubtful, alternative critical-value procedures should be evaluated rather than assuming that orthogonal nuisance estimation resolves the approximation problem. Such reporting does not cure weak relevance, but it separates evidence about validity, numerical scope, and discrimination.

These boundaries define the scope of the conclusions. Within the preserved DGP and implementation, weaker relevance produces larger point-estimation errors, less informative accepted sets, and lower power, with material quantile heterogeneity and imperfect finite-sample calibration in some cells. The practical implication is to report instrument information, point accuracy, coverage, and inferential informativeness separately rather than treating robust test inversion or DML as evidence of precision. Future work could compare alternative DGPs and instrument structures, larger samples, cross-fitted or other nuisance estimators, finite-sample critical values, and empirical applications.
